\subsection{Problem Statement}

Symbolic systems are exact identity capable. Variables denote specific objects, pointers target precise memory locations, and database keys refer to singular records. This capability is not incidental; it is fundamental to symbolic computation, which requires that distinct entities remain distinguishable whenever operations depend on their identity.

Neural embeddings operate differently. They generalize by compressing away semantic detail, mapping distinct inputs to the same representation when those inputs are deemed functionally equivalent. The benefit is generalization: the embedding captures similarity structure and supports efficient nearest-neighbor operations. The cost is collision ambiguity: multiple distinct entities can share the same embedding value, and the representation alone no longer identifies which entity is meant \cite{saunshi2019theoretical, higgins2017beta, alemi2017deep, jing2022understanding, wang2020understanding}.

This paper asks: given a fixed embedding that has already been deployed, what is the minimum price of recovering precise identity?

Formally, let $C$ be a finite class label, let $\pi$ be the representation map, and let $U = \pi(C)$ be the embedding value observed by the decoder. Exact identity recovery asks how much additional description is needed to recover $C$ from $U$. Once the problem is written this way, the central combinatorial quantity is immediate:
\[
A_{\pi} := \max_u |\pi^{-1}(u)|,
\]
the size of the largest collision fiber.

\begin{remark}[Noninjectivity and collision fibers]
A representation $\pi$ is noninjective exactly when at least one fiber $\pi^{-1}(u)$ has size $>1$. The collision multiplicity $A_\pi$ therefore measures the degree of noninjectivity: $A_\pi=1$ iff $\pi$ is injective, and larger $A_\pi$ corresponds to worse collision.
\end{remark}

The worst-case fixed-length law is then exactly the minimum number of extra bits needed to separate that worst fiber:
\[
L \ge \log_2 A_{\pi}.
\]

This viewpoint yields an adaptive law $\lceil \log_2 |\pi^{-1}(u)|\rceil$ per fiber, the fiberwise decomposition of recoverable mass for arbitrary finite sources, and a query-side counting floor $\lceil \log_2 A_{\pi}\rceil \le d$ (since $d$ binary observables induce at most $2^d$ transcripts). Worst-case bits, adaptive bits, query cost, and distortion are different currencies for the same residual ambiguity.

We study the finite, fixed-representation, zero-error regime: the decoder uses only the declared representation $U$ and charged auxiliary description. No hidden registry or external side channel may contribute uncharged distinguishing information. Higher-level regimes must reduce to this base case at deployment time, so they cannot escape the operational implications derived here.

\paragraph{Systems consequence} Symbolic handles are the standard mechanisms for paying the identity debt left by non-injective semantic abstraction. This debt equals the residual ambiguity quantified by collision fiber geometry (Theorems~\ref*{thm:converse} and~\ref*{thm:adaptive-fiber-budget}), resolved through auxiliary bits, query effort, or accepted distortion.

\paragraph{Privacy and security perspective} The fiber geometry directly governs identity disclosure. When distinct entities share the same representation, recovering which specific entity is meant constitutes an \emph{identity disclosure}, and the auxiliary description required to resolve the collision fiber is exactly the information that is protected if privacy is to be maintained. This connects to the Special Issue theme of ``data compression and storage under privacy and security constraints'': our results characterize the precise disclosure-recovery tradeoff for semantic representations. The formal development establishes a finite explicit disclosure separation between nominal-tag access and attribute-only identification, aligning with formal exactness requirements in logic-based explanation frameworks \cite{ignatiev2019abduction, marquessilva2022trustworthy}, complementing broader privacy frameworks such as differential privacy \cite{dwork2006differential}. \leanmeta{\LH{PRIV3}, \LH{PRV1}}

\paragraph{Related frameworks}
The problem relates to Wyner--Ziv side information \cite{slepian1973noiseless, wyner1976rate, yamamoto1981source, oohama1997gaussian, gastpar2004wyner}, zero-error graph-capacity theory \cite{shannon1956zero, lovasz1979shannon}, and semantics-aware communication \cite{gunduz2023beyond, uysal2022semantic, bao2011towards, kountouris2021semantics}. In the fixed-representation regime, our confusability graphs are disjoint cliques, yielding exact finite formulas. The fiber geometry also connects to information bottleneck \cite{Tishby2000} and rate-distortion-perception \cite{blau2019rethinking, chen2022rdp, xie2025output}.

\paragraph{Graph-theoretic interpretation}
The strong product structure of confusability graphs under repeated composition connects this work to the Shannon capacity program \cite{simonyi1995graph, alon1996source, haemers1978upper, zuiddam2019asymptotic} and zero-error systems theory \cite{witsenhausen1976zero}.

\subsection{Contributions}

The collision fiber geometry governs exact identity-recovery costs. Three interlocking threads:

\begin{enumerate}
\item \textbf{Rate-distortion and exact coding.} Exact fiberwise rate-distortion law for arbitrary finite sources via recoverable-mass decomposition, with closed-form $D^\star(L)=\max(0,1-2^L/a)$ for uniform single-block sources. Supporting: tight converse $L \ge \log_2 A_{\pi}$, exact finite-block scaling in the finite-regime spirit of \cite{polyanskiy2010channel}, and adaptive budget $\lceil \log_2 |\pi^{-1}(u)|\rceil$ per fiber.
\item \textbf{Query currency and canonical structure.} Complete derivability-aware families reduce to an orthogonal semantically minimal core, where minimal distinguishing families form matroid bases and query cost obeys $\lceil \log_2 A_{\pi}\rceil \le d$.
\item \textbf{Open-world robustness.} Collision-freedom does not persist under extension; barrier-freedom certification is uncomputable for arbitrary future-generating code.
\end{enumerate}

Remaining payoffs: symbolic handles are necessary in semantics-aware systems; Lean~4 verification provides mechanized integrity.

\paragraph{Machine-checked proofs}
All principal theorems are formalized in Lean~4; informal arguments were admitted only after formal proofs compiled without \texttt{sorry}. Superscript Lean codes reference specific theorems. \leanmeta{\LH{PRV1}}

\paragraph{Worked example}

We use one small example throughout the paper. Five classes have profiles
\[
A=000,\quad B=000,\quad C=100,\quad D=110,\quad E=111.
\]
Then $A$ and $B$ form a collision fiber of size two, so the exact worst-case fixed-length budget is one bit. The adaptive picture is equally simple: the fiber over profile $000$ needs one auxiliary bit, while the singleton fibers need none. Under the adaptive scheme, the encoder sends one bit only when $U=000$ and sends no extra bit for $U\in\{100,110,111\}$.

\subsection{Paper Organization}

The paper develops three interlocking threads. Sections~\ref{sec:model} and~\ref{sec:bits} establish the bit-currency backbone. Section~\ref{sec:queries} develops the query-currency thread and its canonical matroid structure. Section~\ref{sec:distortion} gives the distortion currency via fiberwise recoverable-mass decomposition. Section~\ref{sec:systems} explains why symbolic identity layers are necessary in semantics-aware systems. Section~\ref{sec:robustness} records open-world consequences of the fiber-geometry framework, and Section~\ref{sec:conclusion} closes the paper.
